\section{Introduction}

Backpropagation (BP) gives every parameter, at every step, an error signal
that is globally informed: it reflects the loss, the targets, and the state of
every layer above. This is expensive in compute (a full backward pass costs
roughly twice the forward pass), in memory (activations must be stored to
full depth), and in biological plausibility (exact weight transport and a
global synchronized backward phase have no accepted neural analogue). A long
line of work attacks each cost separately. We instead ask a deliberately
simple question: \emph{along which axes, and by how much, can the BP signal be
degraded before learning breaks?}

We study two independent axes of degradation in small, fully controlled
settings.

\paragraph{Axis 1: locality of the learning signal (what each synapse may
know).} Supermask results show that a randomly weighted, sufficiently wide
network already contains subnetworks that solve the task; learning can be
reduced to \emph{connectivity search}, with weights never updated
\citep{ramanujan2020whats, zhou2019deconstructing, malach2020proving}. But the
known search procedure (edge-popup) uses full backpropagation on mask scores,
so the demonstration that ``learning is wiring'' still smuggles in a global
gradient. We hold the setting fixed---frozen random MLP, learnable binary
mask---and vary only the information available to the mask update rule:
(a)~the full gradient (\ep{}, the upper bound), (b)~a fixed random projection
of the output error vector delivered directly to each layer
(\dfa{} \citep{nokland2016direct}, applied here to mask scores rather than
weights), and (c)~a single global scalar per sample modulating a Hebbian
eligibility trace (\rmhebb{}, a three-factor rule
\citep{fremaux2016neuromodulated, williams1992simple}).

\paragraph{Axis 2: frequency of deep credit assignment (when the network pays
for depth).} Late layers can be updated by backpropagating through only the
final two weight matrices; the trunk below is touched only by full-depth
steps. We propose \fw{} (\emph{front wheeler}): run cheap last-two-layer
updates by default, and fire a \emph{burst} of full-BP steps only when a loss
plateau indicates the cheap subspace is exhausted. An exponential backoff
handles the ``boy who cried wolf'' failure: if a burst itself yields no
improvement, the plateau is global rather than the front's fault, and the
evidence threshold for the next alarm is raised.

\paragraph{Contributions.} On MNIST \citep{lecun1998gradient} and CIFAR-10
\citep{krizhevsky2009learning}, with MLPs and three seeds throughout:
\begin{enumerate}
  \item \textbf{Wiring vs.\ weights.} Gradient mask search over frozen random
  weights comes within half a point of dense training on MNIST ($98.2\%$ vs.\
  $98.7\%$) and \emph{beats} it on CIFAR-10 ($55.9\%$ vs.\ $53.0\%$), where
  dense MLP training overfits (Table~\ref{tab:exp1}).
  \item \textbf{An information threshold for local wiring rules.} A random
  projection of the error \emph{vector} suffices: \dfa{} mask search reaches
  within $1.3$--$2.6$ points of \ep{} and matches dense weight training on
  CIFAR-10. The global-\emph{scalar} rule we stabilized does not suffice:
  \rmhebb{} learns ($2.6$--$5\times$ chance) but plateaus far below all
  vector-signal rules. In these experiments, the threshold for effective
  wiring search sits between a broadcast vector and a broadcast scalar.
  \item \textbf{Adaptive beats scheduled deep credit assignment.} At a matched
  budget of ${\sim}4$--$5\%$ full-BP steps, \fw{} beats the matched periodic
  schedule by $+1.5$ (MNIST) and $+5.1$ (CIFAR-10) points, and beats
  one-switch linear-probe-then-fine-tune \citep{kumar2022finetuning} by
  $+2.0$ and $+6.1$ points respectively (Table~\ref{tab:exp2}). The margins
  grow with task difficulty. A block-coordinate ablation (\bcd{}: bursts
  update the trunk only) shows the closed-loop trigger, not the block
  structure, carries most of the benefit.
  \item \textbf{Two documented failure modes.} (a)~Selecting mask entries by
  $|s|$ while passing straight-through gradients to $s$ silently drops a
  $\operatorname{sign}(s)$ factor and reduces learning to chance.
  (b)~\emph{Isolated} full-BP steps interleaved with front training can
  collapse the network: the front-trained, confident head back-propagates
  error signals into the stale trunk whose norm we measure growing to
  ${\sim}20\times$ its initial scale, driving dead-ReLU collapse unless
  gradients are clipped. The latter is an independent, from-scratch manifestation of the
  feature-distortion mechanism that motivates \lpft{}
  \citep{kumar2022finetuning}.
\end{enumerate}

We emphasize scope: all results are on small MLPs and standard image
classification toys. The contribution is a controlled characterization of two
axes of BP removal---not a claim about state-of-the-art training. Code and
all run logs are released.
